\documentclass[11pt,a4paper]{article}
\usepackage[utf8]{inputenc}
\usepackage[T1]{fontenc}
\usepackage[french]{babel}
\usepackage{geometry}
\usepackage{booktabs}
\usepackage{longtable}
\usepackage{array}
\usepackage{xcolor}
\usepackage{colortbl}
\usepackage{graphicx}
\usepackage{fancyhdr}
\usepackage{titlesec}
\usepackage{tikz}
\usepackage{enumitem}
\usepackage{hyperref}
\usepackage{multicol}
\usepackage{listings}
\usepackage{tcolorbox}
\usepackage{amssymb}


\geometry{margin=2cm}
\pagestyle{fancy}
\fancyhf{}
\rhead{Atelier Machine Learning}
\lhead{Analyse Comportementale Clientèle}
\rfoot{Page \thepage}

% Couleurs personnalisées
\definecolor{headerblue}{RGB}{41, 65, 114}
\definecolor{lightgray}{RGB}{245, 245, 245}
\definecolor{accentorange}{RGB}{230, 126, 34}
\definecolor{codegreen}{RGB}{0,128,0}
\definecolor{codegray}{RGB}{128,128,128}
\definecolor{codepurple}{RGB}{128,0,128}
\definecolor{backcolour}{RGB}{248,248,248}
\definecolor{warningred}{RGB}{231, 76, 60}
\definecolor{successgreen}{RGB}{46, 204, 113}

\lstdefinestyle{mystyle}{
    backgroundcolor=\color{backcolour},
    commentstyle=\color{codegreen},
    keywordstyle=\color{codepurple},
    numberstyle=\tiny\color{codegray},
    stringstyle=\color{codegreen},
    basicstyle=\ttfamily\footnotesize,
    breakatwhitespace=false,
    breaklines=true,
    captionpos=b,
    keepspaces=true,
    numbers=left,
    numbersep=5pt,
    showspaces=false,
    showstringspaces=false,
    showtabs=false,
    tabsize=2
}
\lstset{style=mystyle}

\tcbset{
    colback=backcolour,
    colframe=headerblue,
    fonttitle=\bfseries,
    boxrule=1pt,
    arc=3mm
}

\titleformat{\section}
{\color{headerblue}\normalfont\Large\bfseries}
{\thesection}{1em}{}[{\titlerule[2pt]}]

\titleformat{\subsection}
{\color{accentorange}\normalfont\large\bfseries}
{\thesubsection}{1em}{}

% Définition de colonnes personnalisées
\newcolumntype{C}[1]{>{\centering\arraybackslash}p{#1}}
\newcolumntype{L}[1]{>{\raggedright\arraybackslash}p{#1}}

\begin{document}

% ─────────────────────────────────────────────────────────────
%  PAGE DE GARDE
% ─────────────────────────────────────────────────────────────
\begin{titlepage}
\begin{tikzpicture}[remember picture,overlay]
    \fill[headerblue] (current page.north west) rectangle ([yshift=-5cm]current page.north east);
    \node[anchor=north west, xshift=2cm, yshift=-2cm, text=white] at (current page.north west)
        {\Huge\textbf{Atelier Machine Learning}};
    \node[anchor=north west, xshift=2cm, yshift=-3.5cm, text=white] at (current page.north west)
        {\Large\textit{Analyse Comportementale Clientèle Retail}};
\end{tikzpicture}

\vspace{6cm}

\begin{center}
{\LARGE\textbf{COMPTE RENDU}}\par
\vspace{1cm}
{\large Atelier Pratique : E-commerce de Cadeaux}\par
\vspace{2cm}
{\large \textbf{Préparé par Med Ayoub Ben Ayed}}
\vspace{3cm}

\begin{tikzpicture}
    \draw[headerblue, line width=2pt] (0,0) -- (10,0);
\end{tikzpicture}

\vspace{2cm}

{\textbf{Année Universitaire : 2025--2026}}
\end{center}
\vfill
\end{titlepage}

\tableofcontents
\newpage

% ─────────────────────────────────────────────────────────────
%  1. INTRODUCTION
% ─────────────────────────────────────────────────────────────
\section{Introduction}

Ce compte rendu présente le travail réalisé dans le cadre de l'atelier \textit{Machine Learning appliqué au commerce de détail}. L'objectif principal est de construire un système complet de prédiction du \textbf{churn} (attrition client) à partir de données transactionnelles d'une boutique e-commerce spécialisée dans les cadeaux.

Le projet couvre l'intégralité du cycle de vie d'un projet de machine learning :
\begin{itemize}[noitemsep]
    \item exploration et nettoyage des données,
    \item ingénierie de features et construction d'un pipeline de prétraitement,
    \item entraînement et comparaison de plusieurs modèles de classification,
    \item évaluation des performances,
    \item déploiement d'une API de prédiction via Flask.
\end{itemize}

\begin{tcolorbox}[title=Contexte métier]
Le churn désigne le fait qu'un client cesse d'acheter auprès d'une entreprise. Le détecter à l'avance permet de mettre en place des actions de fidélisation ciblées, réduisant ainsi le coût d'acquisition de nouveaux clients.
\end{tcolorbox}

% ─────────────────────────────────────────────────────────────
%  2. DESCRIPTION DU JEU DE DONNÉES
% ─────────────────────────────────────────────────────────────
\section{Description du jeu de données}

\subsection{Vue d'ensemble}

Le fichier source \texttt{data/raw/retail\_customers\_COMPLETE\_CATEGORICAL.csv} contient les enregistrements de 4~372 clients avec 52 variables (51 features + 1 variable cible).

\begin{center}
\begin{tabular}{lc}
\toprule
\textbf{Caractéristique} & \textbf{Valeur} \\
\midrule
Nombre de clients & 4~372 \\
Nombre de variables & 52 \\
Variable cible & \texttt{Churn} (0 = non-churn, 1 = churn) \\
Taux de churn & 33,3~\% (1~454 / 4~372) \\
\bottomrule
\end{tabular}
\end{center}

\subsection{Variables du dataset}

Les 51 features sont regroupées en plusieurs familles :

\begin{longtable}{L{4.5cm} L{6cm} C{3cm}}
\toprule
\textbf{Variable} & \textbf{Description} & \textbf{Type} \\
\midrule
\endhead
\midrule
\multicolumn{3}{r}{\textit{Suite page suivante}} \\
\endfoot
\bottomrule
\endlastfoot
% RFM
\texttt{Recency} & Nombre de jours depuis le dernier achat (1--374, moy. 92) & Numérique \\
\texttt{Frequency} & Nombre de transactions (1--248, moy. 5) & Numérique \\
\texttt{MonetaryTotal} & Montant total dépensé en £ (moy. 1~898~£) & Numérique \\
\texttt{MonetaryAvg} & Montant moyen par transaction & Numérique \\
\texttt{MonetaryStd} & Écart-type du montant par transaction & Numérique \\
\texttt{MonetaryMin / MonetaryMax} & Montant min/max par transaction & Numérique \\
\midrule
% Transactional
\texttt{TotalQuantity} & Quantité totale de produits commandés & Numérique \\
\texttt{AvgQuantityPerTransaction} & Quantité moyenne par transaction & Numérique \\
\texttt{CancelledTransactions} & Nombre de transactions annulées & Numérique \\
\texttt{ReturnRatio} & Ratio de retours & Numérique \\
\texttt{UniqueInvoices} & Nombre de factures uniques & Numérique \\
\texttt{UniqueProducts} & Nombre de références produits distinctes & Numérique \\
\texttt{AvgLinesPerInvoice} & Lignes moyennes par facture & Numérique \\
\midrule
% Temporel
\texttt{CustomerTenureDays} & Ancienneté du client en jours & Numérique \\
\texttt{RegistrationDate} & Date d'inscription & Date \\
\texttt{PreferredDayOfWeek} & Jour préféré d'achat & Numérique \\
\texttt{PreferredHour} & Heure préférée d'achat & Numérique \\
\texttt{WeekendPurchaseRatio} & Proportion d'achats le week-end & Numérique \\
\texttt{AvgDaysBetweenPurchases} & Intervalle moyen entre achats & Numérique \\
\midrule
% Profil
\texttt{Age} & Âge du client (moy. 49 ans, 30~\% manquants) & Numérique \\
\texttt{Gender} & Genre & Catégorielle \\
\texttt{Country} & Pays (majoritairement UK : 90~\%) & Catégorielle \\
\texttt{Region} & Région géographique & Catégorielle \\
\texttt{AccountStatus} & Statut du compte & Catégorielle \\
\texttt{SupportTicketsCount} & Nombre de tickets de support & Numérique \\
\texttt{SatisfactionScore} & Score de satisfaction (1--5) & Numérique \\
\midrule
% Catégorielles ordonnées
\texttt{SpendingCategory} & Catégorie de dépense (Low/Medium/High/VIP) & Ordinale \\
\texttt{LoyaltyLevel} & Niveau de fidélité (Nouveau/Jeune/Établi/Ancien) & Ordinale \\
\texttt{BasketSizeCategory} & Taille du panier (Petit/Moyen/Grand) & Ordinale \\
\texttt{AgeCategory} & Tranche d'âge & Ordinale \\
\texttt{PreferredTimeOfDay} & Créneau horaire préféré & Ordinale \\
\midrule
% Catégorielles nominales
\texttt{RFMSegment} & Segment RFM (Champions/Fidèles/Potentiels/Dormants) & Nominale \\
\texttt{CustomerType} & Type client (Nouveau/Régulier/Hyperactif/Perdu/Occasionnel) & Nominale \\
\texttt{FavoriteSeason} & Saison de prédilection & Nominale \\
\texttt{ProductDiversity} & Diversité des achats & Nominale \\
\texttt{WeekendPreference} & Préférence semaine/week-end & Nominale \\
\texttt{ChurnRiskCategory} & Catégorie de risque de churn \emph{(variable proxy — supprimée)} & Nominale \\
\end{longtable}

\subsection{Déséquilibre des classes}

Le dataset présente un déséquilibre modéré : 33,3~\% de clients churners contre 66,7~\% de clients actifs. Ce déséquilibre est pris en compte dans le pipeline via la technique SMOTE (voir section~\ref{sec:pipeline}).

\subsection{Valeurs manquantes}

\begin{center}
\begin{tabular}{lcc}
\toprule
\textbf{Variable} & \textbf{Valeurs manquantes} & \textbf{Taux (\%)} \\
\midrule
\texttt{Age} & 1~311 & 30,0 \\
\texttt{AvgDaysBetweenPurchases} & 79 & 1,8 \\
Autres variables & 0 & 0,0 \\
\bottomrule
\end{tabular}
\end{center}

Les valeurs aberrantes détectées dans \texttt{SupportTicketsCount} (valeurs $-1$ et $999$) et \texttt{SatisfactionScore} (valeurs $-1$, $0$ et $99$) sont traitées comme manquantes.

% ─────────────────────────────────────────────────────────────
%  3. EXPLORATION DES DONNÉES
% ─────────────────────────────────────────────────────────────
\section{Exploration des données}

\subsection{Objectifs de l'exploration}

L'exploration a été réalisée dans le notebook \texttt{notebooks/01\_data\_exploration.ipynb}. Elle comprend les étapes suivantes :

\begin{enumerate}[noitemsep]
    \item Inspection de la structure du dataset (\texttt{shape}, \texttt{info()}, \texttt{dtypes}).
    \item Analyse des valeurs manquantes et de leur distribution.
    \item Statistiques descriptives (\texttt{describe()}) sur les variables numériques et catégorielles.
    \item Calcul de la matrice de corrélation et visualisation sous forme de carte thermique.
    \item Filtrage des paires de variables avec une corrélation $|\rho| > 0{,}8$.
    \item Calcul du \textbf{Facteur d'Inflation de la Variance (VIF)} pour détecter la multicolinéarité.
\end{enumerate}

\subsection{Analyse de corrélation}

La matrice de corrélation a révélé de fortes redondances entre certaines variables transactionnelles. Les variables présentant une corrélation supérieure à 0,8 avec une autre variable ont été identifiées pour suppression :

\begin{center}
\begin{tabular}{ll}
\toprule
\textbf{Groupe} & \textbf{Variables à supprimer} \\
\midrule
Montants & \texttt{MonetaryMin}, \texttt{MonetaryMax} \\
Quantités & \texttt{TotalQuantity}, \texttt{MinQuantity}, \texttt{MaxQuantity} \\
Descriptions & \texttt{UniqueDescriptions} \\
Transactions annulées & \texttt{CancelledTransactions} \\
Transactions totales & \texttt{TotalTransactions}, \texttt{UniqueInvoices}, \texttt{AvgLinesPerInvoice} \\
\bottomrule
\end{tabular}
\end{center}

\subsection{Analyse VIF}

L'analyse VIF a identifié trois variables avec un VIF très élevé, indiquant une forte multicolinéarité avec les autres features :

\begin{center}
\begin{tabular}{lc}
\toprule
\textbf{Variable} & \textbf{Action} \\
\midrule
\texttt{CustomerID} & Supprimée (identifiant non prédictif) \\
\texttt{CustomerTenureDays} & Supprimée (VIF élevé) \\
\texttt{FirstPurchaseDaysAgo} & Supprimée (VIF élevé) \\
\bottomrule
\end{tabular}
\end{center}

\subsection{Variables proxy supprimées}

Les variables \texttt{ChurnRiskCategory}, \texttt{NewsletterSubscribed} et \texttt{LastLoginIP} ont été supprimées car elles constituent soit des proxies directs de la variable cible (risque de fuite de données), soit des informations inutilisables en production.

% ─────────────────────────────────────────────────────────────
%  4. PIPELINE DE PRÉTRAITEMENT
% ─────────────────────────────────────────────────────────────
\section{Pipeline de prétraitement}
\label{sec:pipeline}

\subsection{Architecture générale}

Le prétraitement est entièrement encapsulé dans un pipeline scikit-learn/imbalanced-learn défini dans \texttt{src/preprocessing.py}. Cette architecture garantit l'absence de \textit{data leakage} : toutes les transformations (imputation, normalisation, PCA) sont ajustées uniquement sur les données d'entraînement et appliquées telles quelles aux données de test et de production.

\begin{center}
\begin{tikzpicture}[
    box/.style={draw=headerblue, fill=lightgray, rounded corners,
                minimum width=5cm, minimum height=0.75cm,
                text centered, font=\small},
    arrow/.style={->, thick, headerblue}
]
% vertical chain
\node[box] (raw)   at (0,  0)   {Données brutes (52 colonnes)};
\node[box] (clean) at (0, -1.2) {\textbf{DataCleaner} : nettoyage + feature engineering};
\node[box] (dyn)   at (0, -2.4) {\textbf{DynamicColumnTransformer}};
% three branches
\node[box] (num)   at (-4.5, -3.8) {Numériques\\Imputer $\to$ Scaler $\to$ PCA};
\node[box] (ord)   at ( 0,   -3.8) {Ordinales\\Imputer $\to$ OrdinalEncoder};
\node[box] (nom)   at ( 4.5, -3.8) {Nominales\\Imputer $\to$ OneHotEncoder};
% SMOTE + clf
\node[box] (smote) at (0, -5.2) {\textbf{SMOTE} (sur-échantillonnage)};
\node[box] (clf)   at (0, -6.4) {\textbf{Classifieur}};

\draw[arrow] (raw)   -- (clean);
\draw[arrow] (clean) -- (dyn);
\draw[arrow] (dyn)   -- (num);
\draw[arrow] (dyn)   -- (ord);
\draw[arrow] (dyn)   -- (nom);
\draw[arrow] (num)   -- (smote);
\draw[arrow] (ord)   -- (smote);
\draw[arrow] (nom)   -- (smote);
\draw[arrow] (smote) -- (clf);
\end{tikzpicture}
\end{center}

\subsection{Étape 1 : DataCleaner}

La classe \texttt{DataCleaner} (fichier \texttt{src/preprocessing.py}, lignes 88--127) effectue les opérations suivantes :

\begin{enumerate}[noitemsep]
    \item \textbf{Suppression des colonnes} redondantes, identifiants et proxies (22 colonnes supprimées).
    \item \textbf{Correction des valeurs aberrantes} :
        \begin{itemize}[noitemsep]
            \item \texttt{SupportTicketsCount} : $-1$ et $999 \to$ NaN
            \item \texttt{SatisfactionScore} : $-1$, $0$ et $99 \to$ NaN
        \end{itemize}
    \item \textbf{Parsing de la date d'inscription} (\texttt{RegistrationDate}) en quatre variables numériques : \texttt{RegYear}, \texttt{RegMonth}, \texttt{RegDay}, \texttt{RegWeekday}.
    \item \textbf{Ingénierie de features} :
        \begin{itemize}[noitemsep]
            \item $\texttt{MonetaryPerDay} = \frac{\texttt{MonetaryTotal}}{\texttt{Recency} + 1}$
            \item $\texttt{AvgBasketValue} = \frac{\texttt{MonetaryTotal}}{\texttt{Frequency}}$
        \end{itemize}
\end{enumerate}

\subsection{Étape 2 : DynamicColumnTransformer}

Ce transformateur détecte automatiquement les colonnes disponibles après nettoyage et les route vers l'un des trois sous-pipelines :

\begin{center}
\begin{tabular}{L{3.5cm} L{5cm} L{5cm}}
\toprule
\textbf{Branche} & \textbf{Colonnes traitées} & \textbf{Transformations} \\
\midrule
Numériques & Toutes les colonnes \texttt{float/int} & Imputation médiane $\to$ StandardScaler $\to$ PCA (95~\% variance) \\
Ordinales & \texttt{SpendingCategory}, \texttt{LoyaltyLevel}, \texttt{BasketSizeCategory}, \texttt{AgeCategory}, \texttt{PreferredTimeOfDay} & Imputation mode $\to$ OrdinalEncoder (ordres définis) \\
Nominales & \texttt{RFMSegment}, \texttt{CustomerType}, \texttt{FavoriteSeason}, \texttt{Region}, \texttt{Gender}, \texttt{Country}, etc. & Imputation mode $\to$ OneHotEncoder \\
\bottomrule
\end{tabular}
\end{center}

La PCA est configurée pour conserver les composantes expliquant au moins 95~\% de la variance totale des variables numériques, réduisant ainsi la dimensionnalité tout en préservant l'information essentielle.

\subsection{Étape 3 : SMOTE}

La technique SMOTE (\textit{Synthetic Minority Over-sampling Technique}) génère des exemples synthétiques de la classe minoritaire (churn = 1) dans l'espace des features transformées. Elle est appliquée \textbf{uniquement pendant l'entraînement}, ce qui garantit qu'aucune donnée synthétique ne contaminent l'ensemble de test.

\begin{lstlisting}[language=Python, caption=Construction du pipeline complet (src/preprocessing.py)]
def create_training_pipeline(..., use_smote=True) -> ImbPipeline:
    preprocessor = create_preprocessor(...)
    steps = list(preprocessor.steps)    # cleaner + dynamic transformer
    if use_smote:
        steps.append(("smote", SMOTE(random_state=42)))
    steps.append(("classifier", classifier))
    return ImbPipeline(steps=steps)
\end{lstlisting}

% ─────────────────────────────────────────────────────────────
%  5. MODÈLES ENTRAÎNÉS
% ─────────────────────────────────────────────────────────────
\section{Modèles entraînés et comparaison}

\subsection{Séparation train/test}

La séparation des données est réalisée avec une proportion 80~\%/20~\% en utilisant une stratification sur la variable cible afin de conserver le même ratio de churn dans les deux ensembles :

\begin{lstlisting}[language=Python, caption=Séparation stratifiée (src/train\_model.py)]
X_train, X_test, y_train, y_test = train_test_split(
    X, y, test_size=0.2, random_state=42, stratify=y
)
\end{lstlisting}

Cela correspond à environ \textbf{3~497 exemples} d'entraînement et \textbf{875 exemples} de test.

\subsection{Candidats comparés}

Quatre classifieurs sont entraînés et comparés dans \texttt{src/train\_model.py} :

\begin{center}
\begin{tabular}{L{3.5cm} L{8cm} C{2cm}}
\toprule
\textbf{Modèle} & \textbf{Hyperparamètres principaux} & \textbf{Bibliothèque} \\
\midrule
Régression Logistique & \texttt{max\_iter=2000}, \texttt{class\_weight=balanced} & scikit-learn \\
Random Forest & \texttt{n\_estimators=300}, \texttt{class\_weight=balanced} & scikit-learn \\
K-Nearest Neighbors & \texttt{n\_neighbors=9} & scikit-learn \\
XGBoost & \texttt{n\_estimators=300}, \texttt{max\_depth=6}, \texttt{lr=0.05}, \texttt{subsample=0.9} & xgboost \\
\bottomrule
\end{tabular}
\end{center}

\subsection{Critère de sélection}

Le meilleur modèle est sélectionné sur la base du \textbf{score F1} calculé sur l'ensemble de test. Le score F1 est privilégié car il représente la moyenne harmonique entre précision et rappel, ce qui est particulièrement pertinent pour des données déséquilibrées.

$$F_1 = 2 \cdot \frac{\mathrm{Precision} \times \mathrm{Rappel}}{\mathrm{Precision} + \mathrm{Rappel}}$$

\subsection{Métriques d'évaluation}

Les métriques calculées pour chaque modèle sont les suivantes :

\begin{center}
\begin{tabular}{L{4.5cm} L{8cm}}
\toprule
\textbf{Métrique} & \textbf{Description} \\
\midrule
Accuracy & Taux global de bonne classification \\
Précision & Proportion de vrais positifs parmi les prédits positifs \\
Rappel & Proportion de vrais positifs parmi les réels positifs \\
F1-score & Moyenne harmonique précision/rappel \\
ROC-AUC & Aire sous la courbe ROC \\
Matrice de confusion & Distribution des 4 classes (TP, FP, TN, FN) \\
Nombre de composantes PCA & Nombre de composantes retenues pour 95~\% de variance \\
\bottomrule
\end{tabular}
\end{center}

\subsection{Rapport de métriques}

Les résultats sont enregistrés dans \texttt{reports/model\_metrics.json} au format suivant :

\begin{lstlisting}[language=Python, caption=Structure du rapport de métriques]
{
  "best_model": "xgboost",
  "selection_metric": "f1_score",
  "metrics": {
    "logistic_regression": {
      "accuracy": ..., "precision": ..., "recall": ...,
      "f1_score": ..., "roc_auc": ..., "pca_components": ...
    },
    "random_forest": { ... },
    "knn": { ... },
    "xgboost": { ... }
  }
}
\end{lstlisting}

% ─────────────────────────────────────────────────────────────
%  6. DÉPLOIEMENT — API FLASK
% ─────────────────────────────────────────────────────────────
\section{Déploiement : API Flask}

\subsection{Architecture de l'application}

L'application Flask (\texttt{app/app.py}) expose le modèle entraîné sous forme d'une API REST. Le modèle est chargé en mémoire une seule fois au démarrage via la classe \texttt{ChurnPredictor} (\texttt{src/predict.py}).

\begin{center}
\begin{tikzpicture}[
    box/.style={draw=headerblue, fill=lightgray, rounded corners,
                minimum width=3.5cm, minimum height=0.75cm,
                text centered, font=\small},
    arrow/.style={->, thick, headerblue}
]
\node[box] (client) at (0,   0) {Client HTTP};
\node[box] (flask)  at (4,   0) {Flask API (\texttt{app.py})};
\node[box] (pred)   at (8,   0) {\texttt{ChurnPredictor}};
\node[box] (model)  at (12,  0) {Pipeline joblib};

\draw[arrow] (client) -- node[above, font=\scriptsize]{POST /predict} (flask);
\draw[arrow] (flask)  -- node[above, font=\scriptsize]{dict payload}  (pred);
\draw[arrow] (pred)   -- node[above, font=\scriptsize]{predict()}      (model);
\draw[arrow] (model)  -- ++(0,-0.9) -- ++(-12,0) -- node[left, font=\scriptsize]{proba, classe} (client.south);
\end{tikzpicture}
\end{center}

\subsection{Endpoints}

\begin{center}
\begin{tabular}{L{2cm} L{2.5cm} L{4cm} L{4.5cm}}
\toprule
\textbf{Méthode} & \textbf{Route} & \textbf{Corps de requête} & \textbf{Réponse} \\
\midrule
\texttt{GET} & \texttt{/} & — & \texttt{\{"message": "Retail churn API is running"\}} \\
\texttt{POST} & \texttt{/predict} & JSON avec features client & \texttt{\{"churn\_prediction": 0|1, "churn\_probability": float\}} \\
\bottomrule
\end{tabular}
\end{center}

\subsection{Exemple d'utilisation}

\begin{lstlisting}[language=bash, caption=Exemple de requête curl]
curl -X POST http://127.0.0.1:5000/predict \
  -H "Content-Type: application/json" \
  -d '{
    "Recency": 10,
    "Frequency": 15,
    "MonetaryTotal": 1000,
    "RegistrationDate": "2023-01-10",
    "Gender": "F",
    "Country": "France"
  }'
\end{lstlisting}

\begin{lstlisting}[language=Python, caption=Exemple de réponse JSON]
{
  "churn_prediction": 0,
  "churn_probability": 0.12
}
\end{lstlisting}

\subsection{Configuration}

Le chemin du modèle peut être configuré via la variable d'environnement \texttt{CHURN\_MODEL\_PATH}. En l'absence de cette variable, le chemin par défaut \texttt{models/best\_churn\_pipeline.joblib} est utilisé.

La classe \texttt{ChurnPredictor} gère automatiquement les colonnes manquantes dans le payload en les initialisant à \texttt{None}, ce qui permet d'envoyer uniquement les features disponibles sans erreur.

% ─────────────────────────────────────────────────────────────
%  7. TESTS
% ─────────────────────────────────────────────────────────────
\section{Tests}

\subsection{Stratégie de test}

Les tests unitaires et d'intégration sont définis dans \texttt{tests/test\_ml\_pipeline.py} et utilisent des données synthétiques générées à la volée pour ne pas dépendre du fichier CSV de données réelles.

\subsection{Tests implémentés}

\begin{center}
\begin{tabular}{L{5cm} L{8cm}}
\toprule
\textbf{Test} & \textbf{Description} \\
\midrule
\texttt{test\_pipeline\_applies\_pca\_and\_predicts} & Vérifie que le pipeline s'entraîne correctement, génère des prédictions de la bonne longueur, et que la PCA retient entre 1 et 3 composantes sur les features numériques synthétiques. \\
\texttt{test\_train\_and\_predictor\_end\_to\_end} & Test d'intégration complet : entraînement de tous les modèles, vérification de la sauvegarde du modèle et du rapport JSON, puis prédiction via \texttt{ChurnPredictor} sur un client fictif. \\
\bottomrule
\end{tabular}
\end{center}

\subsection{Exécution des tests}

\begin{lstlisting}[language=bash, caption=Commande d'exécution des tests]
python -m unittest -v tests/test_ml_pipeline.py
\end{lstlisting}

% ─────────────────────────────────────────────────────────────
%  8. STRUCTURE DU PROJET
% ─────────────────────────────────────────────────────────────
\section{Structure du projet}

\begin{tcolorbox}[title=Arborescence du projet]
\begin{verbatim}
projet_ml_retail/
├── data/
│   ├── raw/          ← données brutes (CSV source)
│   ├── processed/    ← (placeholder) données nettoyées
│   └── train_test/   ← (généré) X_train/X_test/y_train/y_test
├── notebooks/
│   └── 01_data_exploration.ipynb   ← EDA, corrélation, VIF
├── src/
│   ├── preprocessing.py   ← DataCleaner + pipeline sklearn
│   ├── train_model.py     ← entraînement + sélection du meilleur modèle
│   ├── predict.py         ← ChurnPredictor (inférence)
│   └── utils.py           ← métriques, sauvegarde/chargement modèle
├── app/
│   └── app.py             ← API Flask (POST /predict)
├── models/                ← (généré) best_churn_pipeline.joblib
├── reports/               ← (généré) model_metrics.json
├── tests/
│   └── test_ml_pipeline.py
├── requirements.txt
└── README.md
\end{verbatim}
\end{tcolorbox}

\subsection{Bibliothèques utilisées}

\begin{center}
\begin{tabular}{L{4cm} L{8.5cm}}
\toprule
\textbf{Bibliothèque} & \textbf{Rôle dans le projet} \\
\midrule
\texttt{pandas}, \texttt{numpy} & Chargement et manipulation des données \\
\texttt{scikit-learn $\geq$ 1.3} & Pipelines, transformateurs, modèles, métriques \\
\texttt{imbalanced-learn $\geq$ 0.11} & SMOTE et \texttt{ImbPipeline} \\
\texttt{xgboost $\geq$ 2.0} & Classifieur XGBoost (optionnel) \\
\texttt{joblib $\geq$ 1.3} & Sérialisation du modèle \\
\texttt{flask $\geq$ 3.0} & API REST de prédiction \\
\texttt{matplotlib}, \texttt{seaborn} & Visualisations dans le notebook \\
\texttt{statsmodels $\geq$ 0.14} & Calcul du VIF dans le notebook \\
\bottomrule
\end{tabular}
\end{center}

% ─────────────────────────────────────────────────────────────
%  9. GUIDE D'UTILISATION
% ─────────────────────────────────────────────────────────────
\section{Guide d'utilisation}

\subsection{Installation de l'environnement}

\begin{lstlisting}[language=bash, caption=Mise en place de l'environnement virtuel]
# Créer et activer l'environnement virtuel
python -m venv venv
source venv/bin/activate        # Linux / Mac
# venv\Scripts\Activate.ps1    # Windows PowerShell

# Installer les dépendances
pip install -r requirements.txt
\end{lstlisting}

\subsection{Entraînement des modèles}

\begin{lstlisting}[language=bash, caption=Commande d'entraînement]
python src/train_model.py \
  --data data/raw/retail_customers_COMPLETE_CATEGORICAL.csv \
  --split-output-dir data/train_test \
  --model-output models/best_churn_pipeline.joblib \
  --report-output reports/model_metrics.json \
  --target Churn \
  --pca-variance 0.95
\end{lstlisting}

\textbf{Sorties générées :}
\begin{itemize}[noitemsep]
    \item \texttt{data/train\_test/X\_train.csv}, \texttt{X\_test.csv}, \texttt{y\_train.csv}, \texttt{y\_test.csv}
    \item \texttt{models/best\_churn\_pipeline.joblib}
    \item \texttt{reports/model\_metrics.json}
\end{itemize}

\subsection{Lancement de l'API}

\begin{lstlisting}[language=bash, caption=Démarrage du serveur Flask]
python app/app.py
# Serveur disponible sur http://127.0.0.1:5000
\end{lstlisting}

\subsection{Prédiction en Python}

\begin{lstlisting}[language=Python, caption=Utilisation de ChurnPredictor]
from src.predict import ChurnPredictor

predictor = ChurnPredictor(model_path="models/best_churn_pipeline.joblib")
payload = {
    "Recency": 85,
    "Frequency": 5,
    "MonetaryTotal": 120.0,
    "Country": "France",
    "Gender": "F",
}
result = predictor.predict(payload)
print(result)
# {"churn_prediction": 1, "churn_probability": 0.78}
\end{lstlisting}

% ─────────────────────────────────────────────────────────────
%  10. CONCLUSION
% ─────────────────────────────────────────────────────────────
\section{Conclusion}

Ce projet a permis de construire un pipeline de machine learning complet et robuste pour la prédiction du churn dans un contexte e-commerce retail. Les principaux apports sont :

\begin{itemize}[noitemsep]
    \item Un pipeline de prétraitement \textbf{leakage-safe} entièrement encapsulé dans scikit-learn, reproductible et déployable directement en production.
    \item Une ingénierie de features s'appuyant sur les insights de l'analyse exploratoire (suppression de variables corrélées, VIF, feature engineering RFM).
    \item La prise en compte du déséquilibre de classes via SMOTE et les poids de classes dans les classifieurs.
    \item Un système de comparaison automatique de quatre modèles (Régression Logistique, Random Forest, KNN, XGBoost) avec sélection du meilleur sur la base du F1-score.
    \item Une API Flask prête à l'emploi permettant d'intégrer le modèle dans n'importe quelle application web ou mobile.
\end{itemize}

\begin{tcolorbox}[title=Pistes d'amélioration]
\begin{enumerate}[noitemsep]
    \item Ajouter une \textbf{validation croisée} (StratifiedKFold) pour obtenir des estimations de métriques plus robustes.
    \item Implémenter une \textbf{recherche d'hyperparamètres} (GridSearchCV / Optuna).
    \item Versionner les modèles avec un outil comme \textbf{MLflow} pour le suivi des expériences.
    \item Ajouter un \textbf{pipeline CI/CD} (GitHub Actions) pour automatiser les tests.
    \item Remplacer le serveur de développement Flask par un serveur WSGI de production (\textbf{Gunicorn}).
    \item Ajouter une \textbf{validation des entrées} (Pydantic) à l'API pour sécuriser les requêtes.
\end{enumerate}
\end{tcolorbox}

\end{document}
